\cotefile{156-9}
\scanpage{134}
\mspage{5}
% archive sheet number: 133

\uline{Définition} (provisoire) \unsure{soient|soit} $A$\unsure{,|.} $A'$ deux ateliers. Un \deleted{\unsure{un}} homomorphisme (logique, sans \unsure{connotation} \unsure{géométrique|géometrique} fixée, pour le moment) de $A$ dans $A'$, est une application \deleted{de $A$}
\[
f:A\longrightarrow\inserted{\widehat{A}':=}A'\amalg\{\varnothing\}\quad(\hookrightarrow\unsure{F|\mathcal{F}}(A'))
\]
satisfaisant les trois conditions, pour $X,Y\in A$
\[
\begin{aligned}
\text{\unsure{i)|a)}}\quad &X\trianglelefteq Y\Longrightarrow f(X)\trianglelefteq f(Y)\\
\text{\unsure{ii)|b)}}\quad &X\overset{\circ}{\ll}Y\Longrightarrow f(X)\overset{\circ}{\ll}f(Y)\\
\text{\unsure{iii)|c)}}\quad &X\mathrel{|\circ|}Y\Longrightarrow f(X)\mathrel{|\circ|}f(Y).
\end{aligned}
\]
% margin note written sideways
% slanted
\margin{Mais la condition \unsure{d)} de \unsure{GF|CF} $\overline{\text{\uline{VIII}}}$ \unsure{p.~27|n'est|n'en} \unsure{n'est|vaut|tient|valable} \unsure{pas|pour} \unsure{compte|comptée} pour \unsure{le} moment ! \unsure{!)|i)}}

Pour donner un sens à cette définition, faut donner un sens aux relations
\[
X'\trianglelefteq Y',\quad X'\overset{\circ}{\ll}Y',\quad X'\mathrel{|\circ|}Y'\qquad\text{pour }X',Y'\in A'\amalg\{\varnothing\},
\]
\unsure{des symboles}, quand l'un ou l'autre des deux termes est \deleted{$\{\varnothing\}$} égal à $\varnothing$. Les bonnes conventions sont pour $\forall\,X'\in A'$
\[
\text{a)}\quad X'\trianglelefteq\varnothing,\quad
\deleted{\unsure{\forall X'}}\quad
(\varnothing\trianglelefteq X'\ \deleted{\unsure{\Longleftrightarrow}}\ \text{tjs vrai}
\]
\margin{(\uline{NB} $\trianglelefteq$ sur $\widehat{A}'\deleted{\unsure{\amalg\{\varnothing\}\subset A}}$ n'est \unsure{pas} une relation d'ordre !)}
\[
\text{b)}\quad\varnothing\overset{\circ}{\ll}X',\qquad
(X'\overset{\circ}{\ll}\varnothing\Longleftrightarrow X'=\varnothing)
\]
\margin{\uline{NB} $\overset{\circ}{\ll}$ sur $\widehat{A}'$ est une rel. d'ordre admettant $\varnothing$ comme plus petit él.}
